%% kapa: exam=1 round=29 year=2018 period=1 form=A subject=civil_law subject_ko=민법 seq=35 number=35
%% source: Q-Net 감정평가사 기출문제 (cst00309) artlSeq=5202808 file=extracted/1008/5202808/2018년 제29회 감정평가사 1차시험 기출문제/2018년 제29회 감정평가사 1차시험 기출문제/1교시 A형.pdf
%% conversion: kordoc 4.24.51 -> tools/gen1.py (auto); review=proofread-sonnet; flags=none
\begin{question}{35}
부동산 점유취득시효에 관한 설명으로 옳지 않은 것은? (다툼이 있으면 판례에 따름)
\begin{choices}
  \item 취득시효가 완성되었으나 아직 소유권이전등기를 경료하지 않은 시효완성자에 대하여 소유자는 점유로 인한 부당이득반환청구를 할 수 없다.
  \item 시효기간 중 목적부동산이 제3자에게 양도되어 등기가 이전된 경우, 시효기간 만료시 그 양수인을 상대로 시효취득을 주장할 수 있다.
  \item 소유자가 시효완성 사실을 알고 목적부동산을 제3자에게 처분하고 소유권이전등기를 넘겨준 경우, 소유자는 시효완성자에게 불법행위로 인한 손해배상책임을 진다.
  \item 시효완성자는 취득시효의 기산점과 관련하여 점유기간을 통틀어 등기명의인이 동일한 경우에는 임의의 시점을 기산점으로 할 수 있다.
  \item 소유자가 시효완성 사실을 알고 목적부동산을 제3자에게 처분한 경우, 소유자는 시효완성자에게 채무불이행으로 인한 손해배상책임을 진다.
\end{choices}
\end{question}
\answer{5}
